% Canonical LaTeX source; imported and revised from the legacy Markdown.
\chapter{Goniometria e funzioni iperboliche}
\label{chap:4}

\begin{obiettivocapitolo}{semplificare e risolvere problemi goniometrici con tutte le soluzioni}
\prioritaalta per i prerequisiti di limiti, derivate e integrali.
\end{obiettivocapitolo}
\begin{sceltametodo}{equazioni goniometriche}
\begin{tabularx}{\textwidth}{@{}p{.31\textwidth}X@{}}
$\sin x=a$, $\cos x=a$, $\tan x=a$ & usa la formula generale e poi restringi all'intervallo richiesto.\\
Più funzioni dello stesso angolo & porta tutto a seno e coseno oppure fattorizza.\\
Angoli multipli & usa formule di duplicazione/bisezione o poni $t=\tan(x/2)$ se appropriato.\\
Espressione $A\sin x+B\cos x$ & riscrivi come $R\sin(x+\varphi)$ o $R\cos(x-\varphi)$.\\
Disequazione & risolvi sul periodo fondamentale e replica per periodicità.
\end{tabularx}
\end{sceltametodo}
\begin{proceduraesame}{triangolo generico}
\begin{enumerate}
\item Associa correttamente lati e angoli opposti.
\item Usa il teorema del coseno con tre lati o due lati e angolo compreso; usa il teorema dei seni con coppie lato--angolo note.
\item Nel caso LLA, controlla l'eventuale doppia soluzione prima di scegliere l'arcoseno.
\item Verifica somma degli angoli, positività dei lati e disuguaglianza triangolare.
\end{enumerate}
\end{proceduraesame}
\begin{errorefrequente}{inverse goniometriche}
$\arcsin(\sin x)$ non vale sempre $x$: vale il rappresentante nell'intervallo principale. Riduci l'angolo o usa una definizione a tratti.
\end{errorefrequente}

\section{Angoli, radianti e circonferenza goniometrica}\label{angoli-radianti-e-circonferenza-goniometrica}

\[
180^\circ=\pi\ \mathrm{rad},
\qquad
\theta_{\mathrm{rad}}=\theta_{^\circ}\frac\pi{180},
\qquad
\theta_{^\circ}=\theta_{\mathrm{rad}}\frac{180}{\pi}.
\]

Per un angolo al centro \(\theta\) espresso in radianti, \(r\) è il raggio, \(s\ge0\) la lunghezza dell'arco e \(A_{\mathrm{settore}}\ge0\) l'area del settore:

\[
s=r|\theta|,
\qquad
A_{\mathrm{settore}}=\frac12r^2|\theta|,
\qquad r>0.
\]

Forme inverse:

\[
|\theta|=\frac sr\quad(r>0),
\qquad
r=\frac{s}{|\theta|}\quad(s>0,\ \theta\ne0),
\]

\[
|\theta|=\frac{2A_{\mathrm{settore}}}{r^2}\quad(r>0),
\qquad
r=\sqrt{\frac{2A_{\mathrm{settore}}}{|\theta|}}
\quad(A_{\mathrm{settore}}>0,\ \theta\ne0).
\]

Per \(\theta=0\) si ha \(s=A_{\mathrm{settore}}=0\) per ogni \(r>0\): dai soli valori nulli di arco o area non si può quindi ricavare un raggio unico.

\[
P=(x,y)=(\cos\theta,\sin\theta),
\]

\[
\tan\theta=\frac{\sin\theta}{\cos\theta}
\quad(\cos\theta\ne0),
\qquad
\cot\theta=\frac{\cos\theta}{\sin\theta}
\quad(\sin\theta\ne0),
\]

\[
\sec\theta=\frac1{\cos\theta}
\quad(\cos\theta\ne0),
\qquad
\csc\theta=\frac1{\sin\theta}
\quad(\sin\theta\ne0).
\]

Approfondimento: \href{https://ingegnerismo.it/matematica/angolo/}{angolo}, \href{https://ingegnerismo.it/matematica/cerchio-unitario/}{cerchio unitario}, \href{https://ingegnerismo.it/matematica/trigonometria/}{trigonometria}.

\section{Funzioni goniometriche: domini, immagini, periodicità e valori}\label{funzioni-goniometriche-domini-immagini-periodicituxe0-e-valori}

\[
\begin{array}{c|c|c|c|c}
f&\operatorname{Dom}f&\operatorname{Im}f&T&\text{zeri}\\
\hline
\sin x&\mathbb R&[-1,1]&2\pi&k\pi\\
\cos x&\mathbb R&[-1,1]&2\pi&\pi/2+k\pi\\
\tan x&\mathbb R\setminus\{\pi/2+k\pi\}&\mathbb R&\pi&k\pi\\
\cot x&\mathbb R\setminus\{k\pi\}&\mathbb R&\pi&\pi/2+k\pi\\
\sec x&\mathbb R\setminus\{\pi/2+k\pi\}&(-\infty,-1]\cup[1,+\infty)&2\pi&\varnothing\\
\csc x&\mathbb R\setminus\{k\pi\}&(-\infty,-1]\cup[1,+\infty)&2\pi&\varnothing
\end{array}
\]

con \(k\in\mathbb Z\).

\[
\begin{array}{c|ccc}
\text{quadrante}&\sin&\cos&\tan,\cot\\
\hline
\mathrm I&+&+&+\\
\mathrm{II}&+&-&-\\
\mathrm{III}&-&-&+\\
\mathrm{IV}&-&+&-
\end{array}
\]

\[
\begin{array}{c|ccccc}
\theta&0&\pi/6&\pi/4&\pi/3&\pi/2\\
\hline
\sin\theta&0&1/2&\sqrt2/2&\sqrt3/2&1\\
\cos\theta&1&\sqrt3/2&\sqrt2/2&1/2&0\\
\tan\theta&0&1/\sqrt3&1&\sqrt3&\text{n.d.}\\
\cot\theta&\text{n.d.}&\sqrt3&1&1/\sqrt3&0
\end{array}
\]

\[
\sin(-x)=-\sin x,
\quad
\tan(-x)=-\tan x,
\quad
\cot(-x)=-\cot x,
\quad
\csc(-x)=-\csc x,
\]

\[
\cos(-x)=\cos x,
\qquad
\sec(-x)=\sec x.
\]

Approfondimento: \href{https://ingegnerismo.it/matematica/funzioni-trigonometriche/}{funzioni trigonometriche}, \href{https://ingegnerismo.it/matematica/cerchio-unitario/}{cerchio unitario}.

\section{Angoli associati e identità fondamentali}\label{angoli-associati-e-identituxe0-fondamentali}

\[
\sin^2x+\cos^2x=1,
\]

\[
1+\tan^2x=\sec^2x,
\qquad
1+\cot^2x=\csc^2x,
\]

nei rispettivi domini.

\[
\begin{array}{c|ccc}
\theta&\sin\theta&\cos\theta&\tan\theta\\
\hline
-x&-\sin x&\cos x&-\tan x\\
\pi-x&\sin x&-\cos x&-\tan x\\
\pi+x&-\sin x&-\cos x&\tan x\\
2\pi-x&-\sin x&\cos x&-\tan x\\
\pi/2-x&\cos x&\sin x&\cot x\\
\pi/2+x&\cos x&-\sin x&-\cot x
\end{array}
\]

Approfondimento: \href{https://ingegnerismo.it/matematica/identita-goniometriche/}{identità goniometriche}.

\section{Addizione, sottrazione, duplicazione e bisezione}\label{addizione-sottrazione-duplicazione-e-bisezione}

\[
\sin(a\pm b)=\sin a\cos b\pm\cos a\sin b,
\]

\[
\cos(a\pm b)=\cos a\cos b\mp\sin a\sin b,
\]

\[
\tan(a\pm b)=\frac{\tan a\pm\tan b}{1\mp\tan a\tan b},
\]

nei punti in cui tutti i membri sono definiti.

\[
\cot(a+b)=\frac{\cot a\cot b-1}{\cot a+\cot b},
\qquad
\cot(a-b)=\frac{\cot a\cot b+1}{\cot b-\cot a},
\]

nei punti in cui \(\cot a\), \(\cot b\) e i due membri della formula considerata sono definiti.

\[
\sin2x=2\sin x\cos x,
\]

\[
\cos2x=\cos^2x-\sin^2x=1-2\sin^2x=2\cos^2x-1,
\]

\[
\tan2x=\frac{2\tan x}{1-\tan^2x},
\]

nei rispettivi domini.

\[
\sin3x=3\sin x-4\sin^3x,
\qquad
\cos3x=4\cos^3x-3\cos x.
\]

\[
\sin^2\frac x2=\frac{1-\cos x}{2},
\qquad
\cos^2\frac x2=\frac{1+\cos x}{2},
\]

\[
\sin\frac x2=\pm\sqrt{\frac{1-\cos x}{2}},
\qquad
\cos\frac x2=\pm\sqrt{\frac{1+\cos x}{2}},
\]

con segno determinato dal quadrante di \(x/2\).

\[
\tan\frac x2
=\frac{\sin x}{1+\cos x}
=\frac{1-\cos x}{\sin x},
\]

ciascuna identità nel proprio dominio.

\[
\sin^2x=\frac{1-\cos2x}{2},
\qquad
\cos^2x=\frac{1+\cos2x}{2},
\qquad
\sin x\cos x=\frac12\sin2x.
\]

Approfondimento: \href{https://ingegnerismo.it/matematica/identita-goniometriche/}{identità goniometriche}.

\section{Formule prodotto-somma e somma-prodotto}\label{formule-prodotto-somma-e-somma-prodotto}

\[
\sin a\sin b=\frac12[\cos(a-b)-\cos(a+b)],
\]

\[
\cos a\cos b=\frac12[\cos(a-b)+\cos(a+b)],
\]

\[
\sin a\cos b=\frac12[\sin(a+b)+\sin(a-b)],
\]

\[
\cos a\sin b=\frac12[\sin(a+b)-\sin(a-b)].
\]

\[
\sin a+\sin b=2\sin\frac{a+b}{2}\cos\frac{a-b}{2},
\]

\[
\sin a-\sin b=2\cos\frac{a+b}{2}\sin\frac{a-b}{2},
\]

\[
\cos a+\cos b=2\cos\frac{a+b}{2}\cos\frac{a-b}{2},
\]

\[
\cos a-\cos b=-2\sin\frac{a+b}{2}\sin\frac{a-b}{2}.
\]

Per \((A,B)\ne(0,0)\):

\[
A\sin x+B\cos x=R\sin(x+\varphi),
\qquad
R=\sqrt{A^2+B^2},
\]

\[
\cos\varphi=\frac AR,
\qquad
\sin\varphi=\frac BR,
\qquad
\varphi\pmod{2\pi}.
\]

Per \(A=B=0\), l'espressione è identicamente nulla e \(\varphi\) è indeterminata.

Approfondimento: \href{https://ingegnerismo.it/matematica/identita-goniometriche/}{identità goniometriche}.

\section{Equazioni goniometriche}\label{equazioni-goniometriche}

Le soluzioni sono famiglie infinite. Il parametro \(k\in\mathbb Z\) genera tutti gli angoli coterminali; il simbolo \(\pm\) indica che devono essere considerati entrambi i segni.

Per \(k\in\mathbb Z\):

\[
\sin x=s
\Longleftrightarrow
x=(-1)^k\arcsin s+k\pi,
\qquad |s|\le1,
\]

equivalentemente

\[
x=\arcsin s+2k\pi
\quad\lor\quad
x=\pi-\arcsin s+2k\pi.
\]

\[
\cos x=c
\Longleftrightarrow
x=\pm\arccos c+2k\pi,
\qquad |c|\le1.
\]

\[
\tan x=t
\Longleftrightarrow
x=\arctan t+k\pi,
\]

\[
\cot x=t
\Longleftrightarrow
x=\operatorname{arccot}t+k\pi,
\qquad
\operatorname{arccot}t\in(0,\pi).
\]

\[
\sec x=q
\Longleftrightarrow
\cos x=\frac1q,
\qquad |q|\ge1,
\]

\[
\csc x=q
\Longleftrightarrow
\sin x=\frac1q,
\qquad |q|\ge1.
\]

Per \(R=\sqrt{A^2+B^2}>0\):

\[
A\sin x+B\cos x=C
\Longleftrightarrow
\sin(x+\varphi)=\frac CR.
\]

\[
\begin{array}{c|c}
|C|>R&\varnothing\\
|C|=R&\text{una famiglia modulo }2\pi\\
|C|<R&\text{due famiglie modulo }2\pi
\end{array}
\]

Per \(R=0\): ogni \(x\) se \(C=0\), nessuna soluzione se \(C\ne0\).

Approfondimento: \href{https://ingegnerismo.it/matematica/equazioni-goniometriche-elementari/}{equazioni goniometriche elementari}.

\section{Disequazioni goniometriche}\label{disequazioni-goniometriche}

Per \(k\in\mathbb Z\), \(s,c\in[-1,1]\),

\[
\alpha=\arcsin s\in[-\pi/2,\pi/2],
\qquad
\beta=\arccos c\in[0,\pi].
\]

\[
\sin x\ge s
\Longleftrightarrow
x\in[\alpha+2k\pi,\pi-\alpha+2k\pi],
\]

\[
\sin x\le s
\Longleftrightarrow
x\in[\pi-\alpha+2k\pi,2\pi+\alpha+2k\pi],
\]

\[
\cos x\ge c
\Longleftrightarrow
x\in[-\beta+2k\pi,\beta+2k\pi],
\]

\[
\cos x\le c
\Longleftrightarrow
x\in[\beta+2k\pi,2\pi-\beta+2k\pi].
\]

Per \(t\in\mathbb R\) e \(\gamma=\arctan t\):

\[
\tan x\ge t
\Longleftrightarrow
x\in[\gamma+k\pi,\pi/2+k\pi),
\]

\[
\tan x\le t
\Longleftrightarrow
x\in(-\pi/2+k\pi,\gamma+k\pi].
\]

Per \(\delta=\operatorname{arccot}t\in(0,\pi)\):

\[
\cot x\ge t
\Longleftrightarrow
x\in(k\pi,\delta+k\pi],
\]

\[
\cot x\le t
\Longleftrightarrow
x\in[\delta+k\pi,(k+1)\pi).
\]

Per \(s,c\in(-1,1)\):

\[
\sin x>s
\Longleftrightarrow
x\in(\alpha+2k\pi,\pi-\alpha+2k\pi),
\]

\[
\sin x<s
\Longleftrightarrow
x\in(\pi-\alpha+2k\pi,2\pi+\alpha+2k\pi),
\]

\[
\cos x>c
\Longleftrightarrow
x\in(-\beta+2k\pi,\beta+2k\pi),
\]

\[
\cos x<c
\Longleftrightarrow
x\in(\beta+2k\pi,2\pi-\beta+2k\pi).
\]

\[
\begin{aligned}
\tan x>t
&\Longleftrightarrow
x\in(\gamma+k\pi,\pi/2+k\pi),\\
\tan x<t
&\Longleftrightarrow
x\in(-\pi/2+k\pi,\gamma+k\pi),\\
\cot x>t
&\Longleftrightarrow
x\in(k\pi,\delta+k\pi),\\
\cot x<t
&\Longleftrightarrow
x\in(\delta+k\pi,(k+1)\pi).
\end{aligned}
\]

Per \(h\in\{\sin,\cos\}\):

\[
\begin{array}{c|cccc}
&u<-1&u=-1&u=1&u>1\\
\hline
\{x:h(x)>u\}&\mathbb R&\mathbb R\setminus h^{-1}(\{-1\})&\varnothing&\varnothing\\
\{x:h(x)<u\}&\varnothing&\varnothing&\mathbb R\setminus h^{-1}(\{1\})&\mathbb R\\
\{x:h(x)\ge u\}&\mathbb R&\mathbb R&h^{-1}(\{1\})&\varnothing\\
\{x:h(x)\le u\}&\varnothing&h^{-1}(\{-1\})&\mathbb R&\mathbb R
\end{array}
\]

Approfondimento: \href{https://ingegnerismo.it/matematica/disequazioni-goniometriche/}{disequazioni goniometriche}.

\section{Funzioni goniometriche inverse e composizioni}\label{funzioni-goniometriche-inverse-e-composizioni}

\[
\arcsin:[-1,1]\to[-\pi/2,\pi/2],
\]

\[
\arccos:[-1,1]\to[0,\pi],
\qquad
\arctan:\mathbb R\to(-\pi/2,\pi/2).
\]

\[
\sin(\arcsin y)=y,
\qquad
\cos(\arccos y)=y,
\qquad |y|\le1,
\]

\[
\tan(\arctan y)=y,
\qquad y\in\mathbb R.
\]

\[
\cos(\arcsin y)=\sqrt{1-y^2},
\qquad
\sin(\arccos y)=\sqrt{1-y^2},
\qquad |y|\le1,
\]

\[
\tan(\arcsin y)=\frac{y}{\sqrt{1-y^2}},
\qquad |y|<1,
\]

\[
\tan(\arccos y)=\frac{\sqrt{1-y^2}}{y},
\qquad |y|\le1,
\qquad y\ne0,
\]

\[
\sin(\arctan y)=\frac{y}{\sqrt{1+y^2}},
\qquad
\cos(\arctan y)=\frac1{\sqrt{1+y^2}}.
\]

\[
\arctan(\tan x)=x-k\pi,
\qquad
x\in(-\pi/2+k\pi,\pi/2+k\pi).
\]

\[
\arcsin(\sin x)=
\begin{cases}
x-2k\pi,&x\in[-\pi/2+2k\pi,\pi/2+2k\pi],\\
\pi-x+2k\pi,&x\in[\pi/2+2k\pi,3\pi/2+2k\pi],
\end{cases}
\]

\[
\arccos(\cos x)=
\begin{cases}
x-2k\pi,&x\in[2k\pi,(2k+1)\pi],\\
2k\pi-x,&x\in[(2k-1)\pi,2k\pi].
\end{cases}
\]

Approfondimento: \href{https://ingegnerismo.it/matematica/funzioni-trigonometriche/}{funzioni trigonometriche}, \href{https://ingegnerismo.it/matematica/funzione-inversa/}{funzione inversa}.

\section{Trigonometria dei triangoli}\label{trigonometria-dei-triangoli}

In un triangolo qualunque i lati \(a,b,c\) sono opposti rispettivamente agli angoli \(A,B,C\); \(R\) è il raggio della circonferenza circoscritta, \(K\) l'area e \(s\) il semiperimetro.

\textbf{Triangolo rettangolo.} Se \(C=\pi/2\), allora \(c\) è l'ipotenusa e \(a,b\) sono i cateti:

\[
c^2=a^2+b^2,
\qquad
c=\sqrt{a^2+b^2},
\]

\[
a=\sqrt{c^2-b^2},
\qquad
b=\sqrt{c^2-a^2},
\]

con \(a,b,c>0\) e \(c>a\), \(c>b\).

Rispetto all'angolo acuto \(A\):

\[
\sin A=\frac ac,
\qquad
\cos A=\frac bc,
\qquad
\tan A=\frac ab,
\]

\[
a=c\sin A=b\tan A,
\qquad
b=c\cos A=\frac a{\tan A},
\]

\[
A=\arcsin\frac ac
=\arccos\frac bc
=\arctan\frac ab.
\]

Le formule per l'angolo \(B\) si ottengono scambiando \(a\) e \(b\); inoltre \(A+B=\pi/2\).

\textbf{Triangolo qualunque:}

\[
A+B+C=\pi,
\qquad
A=\pi-B-C,
\]

con \(0<A,B,C<\pi\).

Teorema dei seni:

\[
\frac a{\sin A}=\frac b{\sin B}=\frac c{\sin C}=2R,
\]

\[
a=2R\sin A,
\qquad
R=\frac a{2\sin A},
\qquad
\sin A=\frac a{2R},
\]

\[
a=b\frac{\sin A}{\sin B}
=c\frac{\sin A}{\sin C},
\]

nei casi in cui i denominatori sono non nulli; in un triangolo non degenere lo sono sempre.

Teorema del coseno:

\[
a^2=b^2+c^2-2bc\cos A,
\]

\[
b^2=a^2+c^2-2ac\cos B,
\qquad
c^2=a^2+b^2-2ab\cos C.
\]

Forme inverse per gli angoli:

\[
A=\arccos\frac{b^2+c^2-a^2}{2bc},
\qquad
B=\arccos\frac{a^2+c^2-b^2}{2ac},
\]

\[
C=\arccos\frac{a^2+b^2-c^2}{2ab}.
\]

Condizioni di esistenza del triangolo:

\[
a,b,c>0,
\qquad
a<b+c,
\qquad
b<a+c,
\qquad
c<a+b.
\]

Area:

\[
K=\frac12bc\sin A
=\frac12ca\sin B
=\frac12ab\sin C,
\]

\[
\sin A=\frac{2K}{bc},
\qquad
\sin B=\frac{2K}{ca},
\qquad
\sin C=\frac{2K}{ab},
\]

\[
s=\frac{a+b+c}{2},
\qquad
K=\sqrt{s(s-a)(s-b)(s-c)}.
\]

Scelta della formula in base ai dati:

\begin{longtable}[]{@{}
  >{\raggedright\arraybackslash}p{(\columnwidth - 2\tabcolsep) * \real{0.5000}}
  >{\raggedright\arraybackslash}p{(\columnwidth - 2\tabcolsep) * \real{0.5000}}@{}}
\toprule\noalign{}
\begin{minipage}[b]{\linewidth}\raggedright
Dati noti
\end{minipage} & \begin{minipage}[b]{\linewidth}\raggedright
Formula primaria
\end{minipage} \\
\midrule\noalign{}
\endhead
\bottomrule\noalign{}
\endlastfoot
SSS: tre lati & teorema del coseno inverso \\
SAS: due lati e angolo compreso & teorema del coseno \\
ASA o AAS: due angoli e un lato & somma degli angoli e teorema dei seni \\
SSA: due lati e un angolo non compreso & caso ambiguo del teorema dei seni \\
\end{longtable}

Caso SSA con \(0<A<\pi/2\), lato noto \(a\) opposto ad \(A\), altro lato noto \(b\) e altezza \(h=b\sin A\):

\begin{longtable}[]{@{}lr@{}}
\toprule\noalign{}
Condizione & Numero di triangoli \\
\midrule\noalign{}
\endhead
\bottomrule\noalign{}
\endlastfoot
\(a<h\) & \(0\) \\
\(a=h\) & \(1\), rettangolo \\
\(h<a<b\) & \(2\) \\
\(a\ge b\) & \(1\) \\
\end{longtable}

Per \(A\ge\pi/2\) esiste un solo triangolo se e solo se \(a>b\).

Approfondimento: \href{https://ingegnerismo.it/matematica/trigonometria/}{trigonometria}, \href{https://ingegnerismo.it/matematica/teorema-dei-seni/}{teorema dei seni}, \href{https://ingegnerismo.it/matematica/teorema-del-coseno/}{teorema del coseno}.

\section{Funzioni iperboliche}\label{funzioni-iperboliche}

\[
\sinh x=\frac{e^x-e^{-x}}2,
\qquad
\cosh x=\frac{e^x+e^{-x}}2,
\]

\[
\tanh x=\frac{\sinh x}{\cosh x},
\qquad
\coth x=\frac{\cosh x}{\sinh x}\quad(x\ne0).
\]

\[
\begin{array}{c|c|c|c}
f&\operatorname{Dom}f&\operatorname{Im}f&\text{parità}\\
\hline
\sinh&\mathbb R&\mathbb R&\text{dispari}\\
\cosh&\mathbb R&[1,+\infty)&\text{pari}\\
\tanh&\mathbb R&(-1,1)&\text{dispari}\\
\coth&\mathbb R\setminus\{0\}&(-\infty,-1)\cup(1,+\infty)&\text{dispari}
\end{array}
\]

\[
\cosh^2x-\sinh^2x=1,
\qquad
1-\tanh^2x=\frac1{\cosh^2x}.
\]

\[
\sinh(x\pm y)=\sinh x\cosh y\pm\cosh x\sinh y,
\]

\[
\cosh(x\pm y)=\cosh x\cosh y\pm\sinh x\sinh y.
\]

\[
\sinh2x=2\sinh x\cosh x,
\qquad
\cosh2x=\cosh^2x+\sinh^2x.
\]

\[
\operatorname{arsinh}x=\ln(x+\sqrt{x^2+1}),
\qquad x\in\mathbb R,
\]

\[
\operatorname{arcosh}x=\ln(x+\sqrt{x^2-1}),
\qquad x\ge1,
\]

\[
\operatorname{artanh}x=\frac12\ln\frac{1+x}{1-x},
\qquad |x|<1,
\]

\[
\operatorname{arcoth}x=\frac12\ln\frac{x+1}{x-1},
\qquad |x|>1.
\]

Composizioni diretta dopo inversa, nei rispettivi domini:

\[
\sinh(\operatorname{arsinh}x)=x,
\qquad
\cosh(\operatorname{arcosh}x)=x\quad(x\ge1),
\]

\[
\tanh(\operatorname{artanh}x)=x\quad(|x|<1),
\qquad
\coth(\operatorname{arcoth}x)=x\quad(|x|>1).
\]

Composizioni inversa dopo diretta:

\[
\operatorname{arsinh}(\sinh x)=x,
\qquad
\operatorname{arcosh}(\cosh x)=|x|,
\qquad x\in\mathbb R,
\]

\[
\operatorname{artanh}(\tanh x)=x\quad(x\in\mathbb R),
\qquad
\operatorname{arcoth}(\coth x)=x\quad(x\ne0).
\]

Il valore assoluto nella composizione con \(\cosh\) è necessario perché \(\cosh\) è pari e diventa invertibile soltanto restringendola a \([0,+\infty)\).

Approfondimento: \href{https://ingegnerismo.it/matematica/funzioni-iperboliche/}{funzioni iperboliche}.
